\chapter[Band 18: Rationale Zahlen]{Band 18 auf einen Blick}
\noindent\textbf{Rationale Zahlen}\par\medskip
\noindent\emph{Lesart.}
Bruchpaare besitzen Zähler \(a,c\in\mathbb Z\) und positive Nenner
\(b,d\in\PositiveIntegers\). In rationalen Formeln sind \(p,q,r\in\mathbb Q\).
Die Zeichen \(\RatZero,\RatOne\) unterscheiden rationale Konstanten von
den ganzzahligen. Die Axiomatik des erzeugten geordneten Körpers folgt
im Band auf ihren Modellnachweis.

\begin{BandUebersicht}
\textbf{Bruchpaare und Äquivalenz}
\UebFormel{\begin{aligned}
\RatPairCarrier&\coloneqq
\mathbb Z\times\PositiveIntegers,\\
(a,b)\RatPairRel(c,d)&\coloneqq ad=cb.
\end{aligned}}
\UebQuellen{\FormulaRefAuto{RationalPairCarrier}[def];
\FormulaRefAuto{RationalPairRelDef}[def]}
&
\textit{Äquivalenzrelation auf dem Bruchträger}
\UebFormel{\EqRel{\RatPairCarrier,\RatPairRel}.}
Die Positivität der Nenner erlaubt bei der Transitivität
die Kürzung eines gemeinsamen Faktors.
\UebQuellen{\FormulaRefAuto{RationalPairEquivalence}[thm];
\FormulaRefAuto{RationalPairRelTransitive}[thm]}
\\ \addlinespace[.8ex]

\textbf{Rationale Zahlen und Bruchklassen}
\UebFormel{\begin{aligned}
\mathbb Q&\coloneqq
\QuotientSet{\RatPairCarrier}{\RatPairRel},\\
\RatClass a b&\coloneqq
\EqClass{(a,b)}{\RatPairRel}{\RatPairCarrier}.
\end{aligned}}
\UebQuellen{\FormulaRefAuto{RationalsAsQuotient}[def];
\FormulaRefAuto{RationalClassDef}[def]}
&
\textit{Gleichheit und Darstellung}
\UebFormel{\begin{aligned}
&\RatClass a b=\RatClass c d
\leftrightarrow ad=cb,\\
&q\in\mathbb Q\vdash
\exists a\in\mathbb Z\;\exists b\in\PositiveIntegers\\
&\hspace{5em}q=\RatClass a b.
\end{aligned}}
\UebQuellen{\FormulaRefAuto{RationalClassEquality}[thm];
\FormulaRefAuto{RationalPairRepresentative}[thm]}
\\ \addlinespace[.8ex]

\textbf{Ganzzahlige Einbettung und Zahlkopien}
\UebFormel{\begin{aligned}
&\RatEmbed\colon\mathbb Z\to\mathbb Q,\qquad
\RatEmbed(z)\coloneqq\RatClass z{\IntOne},\\
&\IntCopyInRat\coloneqq\RatEmbed[\mathbb Z],\\
&\NatCopyInRat\coloneqq\RatEmbed[\NatCopyInInt].
\end{aligned}}
\UebQuellen{\FormulaRefAuto{IntegerRationalEmbedding}[def];
\FormulaRefAuto{IntegerCopyInRationals}[def];
\FormulaRefAuto{NaturalCopyInRationals}[def]}
&
\textit{Einbettung des bisherigen Zahlbereichs}
Für \(z,w\in\mathbb Z\):
\UebFormel{\begin{aligned}
&\RatEmbed\colon\mathbb Z\inj\mathbb Q,\\
&\RatEmbed(z+w)=\RatEmbed(z)+\RatEmbed(w),\\
&\RatEmbed(zw)=\RatEmbed(z)\RatEmbed(w),\\
&\NatCopyInRat\subseteq\IntCopyInRat\subseteq\mathbb Q.
\end{aligned}}
\UebQuellen{\FormulaRefAuto{IntegerRationalEmbeddingInjective}[thm];
\FormulaRefAuto{IntegerRationalEmbeddingAdditive}[thm];
\FormulaRefAuto{IntegerRationalEmbeddingMultiplicative}[thm];
\FormulaRefAuto{EmbeddedNaturalIntegerRationalTower}[thm]}
\\ \addlinespace[.8ex]

\textbf{Negation und Addition durch Quotientenabstieg}
Die Definitionen wählen jeweils den eindeutig bestimmten rationalen
Wert aus den Repräsentanten:
\UebFormel{\begin{aligned}
&-q\coloneqq\iota u\in\mathbb Q:\\
&\quad\exists a\in\mathbb Z\;\exists b\in\PositiveIntegers\\
&\quad(q=\RatClass a b\land u=\RatClass{-a}b).
\end{aligned}}
Die Addition verwendet entsprechend den Wert
\(\RatClass{ad+cb}{bd}\) zu \(q=\RatClass a b,\ r=\RatClass c d\).
\UebQuellen{\FormulaRefAuto{RationalNegationDef}[def];
\FormulaRefAuto{RationalAdditionDef}[def]}
&
\textit{Klassenrechnung und additive Gesetze}
\UebFormel{\begin{aligned}
&-\RatClass a b=\RatClass{-a}b,\\
&\RatClass a b+\RatClass c d=\RatClass{ad+cb}{bd},\\
&(p+q)+r=p+(q+r),\\
&q+r=r+q,\qquad q+(-q)=\RatZero.
\end{aligned}}
\UebQuellen{\FormulaRefAuto{RationalNegationClassValue}[thm];
\FormulaRefAuto{RationalAdditionClassValue}[thm];
\FormulaRefAuto{RationalAddAssociative}[thm];
\FormulaRefAuto{RationalAddCommutative}[thm];
\FormulaRefAuto{RationalAddInverse}[thm]}
\\ \addlinespace[.8ex]

\textbf{Multiplikation durch Quotientenabstieg}
\UebFormel{\begin{aligned}
q r\coloneqq\iota u\in\mathbb Q\;&
\exists a,c\in\mathbb Z\\
&\exists b,d\in\PositiveIntegers:\\
&q=\RatClass a b\land r=\RatClass c d\\
&{}\land u=\RatClass{ac}{bd}.
\end{aligned}}
\UebQuellen{\FormulaRefAuto{RationalMultiplicationDef}[def]}
&
\textit{Multiplikative Gesetze}
\UebFormel{\begin{aligned}
&\RatClass a b\cdot\RatClass c d=\RatClass{ac}{bd},\\
&(pq)r=p(qr),\qquad qr=rq,\\
&q\RatOne=q,\qquad p(q+r)=pq+pr.
\end{aligned}}
\UebQuellen{\FormulaRefAuto{RationalMultiplicationClassValue}[thm];
\FormulaRefAuto{RationalMulAssociative}[thm];
\FormulaRefAuto{RationalMulCommutative}[thm];
\FormulaRefAuto{RationalMulOne}[thm];
\FormulaRefAuto{RationalDistributive}[thm]}
\\ \addlinespace[.8ex]

\textbf{Kehrwert und Division}
Für \(a\neq\IntZero\):
\UebFormel{\begin{aligned}
\RatRecipRep a b&\coloneqq
\begin{cases}
\RatClass b a,&\IntZero<a,\\
\RatClass{-b}{-a},&a<\IntZero,
\end{cases}\\
q/r&\coloneqq q\,r^{-1}\quad(r\neq\RatZero).
\end{aligned}}
Der Kehrwert \(q^{-1}\) ist der eindeutige Wert dieser Vorschrift
für einen Repräsentanten von \(q\neq\RatZero\).
\UebQuellen{\FormulaRefAuto{RationalReciprocalRepresentativeDef}[def];
\FormulaRefAuto{RationalReciprocalDef}[def];
\FormulaRefAuto{RationalDivisionDef}[def]}
&
\textit{Wohldefiniertheit und Inversengesetz}
\UebFormel{\begin{aligned}
&\RatClass a b=\RatZero\leftrightarrow a=\IntZero,\\
&q\neq\RatZero\vdash q^{-1}\in\mathbb Q,\\
&q\neq\RatZero\vdash qq^{-1}=\RatOne,\\
&r\neq\RatZero\vdash q/r\in\mathbb Q.
\end{aligned}}
\UebQuellen{\FormulaRefAuto{RationalClassZeroIffNumeratorZero}[thm];
\FormulaRefAuto{RationalReciprocalCompatible}[thm];
\FormulaRefAuto{RationalReciprocalClosure}[thm];
\FormulaRefAuto{RationalMulInverse}[thm];
\FormulaRefAuto{RationalDivisionClosure}[thm]}
\\ \addlinespace[.8ex]

\textbf{Quotientenordnung und strikter Vergleich}
\UebFormel{\begin{aligned}
&\RatOrderGraph\coloneqq\\
&\quad\{(\RatClass a b,\RatClass c d)\mid ad\leq cb\},\\
&q\leq_{\mathbb Q}r\coloneqq(q,r)\in\RatOrderGraph,\\
&q<r\coloneqq q\leq r\land q\neq r.
\end{aligned}}
Die Variablen der Mengenbeschreibung haben die eingangs angegebenen Typen.
\UebQuellen{\FormulaRefAuto{RationalOrderRelationDef}[def];
\FormulaRefAuto{RationalOrderDef}[def];
\FormulaRefAuto{RationalStrictOrderDef}[def]}
&
\textit{Ordnung und Arithmetik sind verträglich}
\UebFormel{\begin{aligned}
&\TotOrd{\mathbb Q,\leq_{\mathbb Q}},\\
&q\leq r\leftrightarrow p+q\leq p+r,\\
&\RatZero<p\land q<r\vdash pq<pr,\\
&z\leq w\leftrightarrow\RatEmbed(z)\leq\RatEmbed(w).
\end{aligned}}
In der letzten Zeile sind \(z,w\in\mathbb Z\).
\UebQuellen{\FormulaRefAuto{RationalOrderTotal}[thm];
\FormulaRefAuto{RationalOrderTranslation}[thm];
\FormulaRefAuto{RationalPositiveMulStrictOrder}[thm];
\FormulaRefAuto{IntegerRationalEmbeddingOrder}[thm]}
\\ \addlinespace[.8ex]

\textbf{Axiomatik des erzeugten geordneten Körpers}
Neben den Körper- und Ordnungsaxiomen wird ausdrücklich
die Erzeugung aus ganzen Zahlen verlangt:
\UebFormel{\begin{aligned}
q\in\mathbb Q\vdash
\exists a\in\mathbb Z\;\exists b\in\PositiveIntegers\\
q=\RatEmbed(a)\bigl(\RatEmbed(b)\bigr)^{-1}.
\end{aligned}}
\UebQuellen{\FormulaRefAuto{RationalOrderedFieldAxioms}[def];
\FormulaRefAuto{RationalFractionGenerationAxiom}[ax]}
&
\textit{Das Quotientenmodell erfüllt die Axiomatik}
\UebFormel{\begin{aligned}
&\mathbb Q=\QuotientSet{\RatPairCarrier}{\RatPairRel},\\
&\RationalOrderedField{\mathbb Q}.
\end{aligned}}
Auch die links als Axiom verwendete Brucherzeugung ist für das
konstruierte Modell ein bewiesenes Resultat.
\UebQuellen{\FormulaRefAuto{RationalFractionGenerationModel}[thm];
\FormulaRefAuto{RationalQuotientModelSatisfiesOrderedField}[thm]}
\\ \addlinespace[.8ex]

\textbf{Positive Nenner und Ganzzahlschranken}
Die axiomatische Weiterarbeit importiert
\UebFormel{\begin{aligned}
&b\in\PositiveIntegers\vdash b\in\mathbb Z
\land\IntZero<b,\\
&a\in\mathbb Z,\ b\in\PositiveIntegers\vdash\\
&\qquad\exists n\in\mathbb N\;a<\IntEmbed(n)b.
\end{aligned}}
\UebQuellen{\FormulaRefAuto{RationalPositiveDenominatorAxiom}[ax];
\FormulaRefAuto{RationalIntegerFractionBoundAxiom}[ax]}
&
\textit{Dichte und Archimedizität}
\UebFormel{\begin{aligned}
&q<r\vdash\exists s\in\mathbb Q\;(q<s<r),\\
&q\in\mathbb Q\vdash\exists n\in\mathbb N\\
&\qquad q<\RatEmbed(\IntEmbed(n)).
\end{aligned}}
Zwischen rationalen Zahlen liegt also wieder eine rationale Zahl;
die natürlichen Bilder sind dennoch nach oben unbeschränkt.
\UebQuellen{\FormulaRefAuto{RationalOrderDense}[thm];
\FormulaRefAuto{RationalArchimedean}[thm]}
\\
\end{BandUebersicht}
